% --- Използвани понятия (незадължителна таблица, по образец) --
\chapter*{Използвани понятия}
\addcontentsline{toc}{chapter}{Използвани понятия}

Таблица~\ref{tab:glossary} съдържа списък на някои от използваните в тази
работа понятия, които не са изрично дефинирани в текста.

\begin{table}[H]
\centering
\renewcommand{\arraystretch}{1.3}
\begin{tabular}{@{}p{0.28\textwidth} p{0.64\textwidth}@{}}
\toprule
\textbf{Понятие} & \textbf{Описание} \\
\midrule
MIDI & Протокол за цифрово представяне на музикални събития (ноти, време, канали, динамика). \\
E-GMD & \eng{Expanded Groove MIDI Dataset} — набор от изпълнения на електронни барабани. \\
Динамика (\eng{velocity}) & Сила/интензитет на удара на нота, кодиран в MIDI като цяло число в диапазона 0--127. \\
Темпо & Скоростта на изпълнение на музикално произведение, обикновено измервана в удари в минута (BPM). \\
Соло барабан (\eng{snare drum}) & Основен барабан в ударната секция, често използван за водещи ритмични фигури и акценти. \\
Каса (\eng{bass drum / kick drum}) & Голям барабан в ударната секция с басов звук, използван за водещи удари предимно на силно време. \\
Фус (\eng{hi-hat}) & Комплект от два чинела, захлюпени един върху друг, които могат да бъдат в отворено (open hi-hat) или затворено положение (closed hi-hat), както и да звучат чрез затваряне с крак (pedal hi-hat). Използвани предимно за поддържане на ритъма. \\
Призрачна нота (\eng{ghost note}) & Много тих удар (обикновено по соло барабан), който запълва груува, без да е самостоятелен акцент. \\
Метрична фаза & Позицията на нотата в рамките на такт или време, изразена като число в $[0,1)$. \\
LightGBM & Табличен модел, използващ градиентен бустинг с дървета, използван тук като базов модел за сравнение. \\
Трансформър (\eng{transformer}) & Архитектура на неврони мрежи, базирана на „механизъм за внимание'' (\eng{attention}). \\
MDN & Мрежа със смес от гаусиани (\eng{mixture density network}) — предсказва многомодално разпределение. \\
Бленда & Конфигурация на електронния комплект барабани, която определя как се обработват входните сигнали и какъв е техният получен звук. \\
Груув (\eng{groove}) & Усешането за ритъм на партията, оформено от вариациите на времето и динамиката. \\
% TODO: добави още понятия при нужда
\bottomrule
\end{tabular}
\caption{Списък на използваните понятия}
\label{tab:glossary}
\end{table}

\clearpage
